\documentclass[11pt]{article}
\usepackage{amsmath,amssymb,amsthm,mathtools}
\usepackage[a4paper,margin=1in]{geometry}
\usepackage{bm}
\usepackage{xcolor}
\usepackage{hyperref}
\newtheorem{theorem}{Theorem}
\newtheorem{remark}[theorem]{Remark}
\title{Discontinuous Galerkin variational formulation for a 1D embedded Blood Flow system with Lax--Friedrichs flux}
\author{}
\date{}

\newcommand{\R}{\mathbb{R}}
\newcommand{\avg}[1]{\left\{\!\!\left\{#1\right\}\!\!\right\}}
\newcommand{\jump}[1]{\left[\!\left[#1\right]\!\right]}
\newcommand{\dt}{\Delta t}
\newcommand{\Vh}{V_h}
\newcommand{\Th}{\mathcal{T}_h}
\newcommand{\Fh}{\mathcal{F}_h}
\newcommand{\Kh}{K}
\newcommand{\norm}[1]{\left\lVert #1\right\rVert}


\begin{document}
\maketitle

\section{Model problem and geometry}
We consider the 1D embedded blood flow model along a centerline
$\Gamma \subset \R^{\text{spacedim}}$:
\begin{align}
  \partial_t A + \bm b\!\cdot\!\nabla (A U) &= f_A,
  && \text{(Area equation)}, \label{eq:area-eq}\\
  \partial_t U + \bm b\!\cdot\!\nabla\!\left(\frac{U^2}{2}\right)
  + \frac{1}{\rho}\bm b\!\cdot\!\nabla P(A)
  + c\,U &= f_U,
  && \text{(Momentum equation)}. \label{eq:momentum-eq}
\end{align}
Here $A$ and $U$ denote the cross-sectional area and mean velocity, respectively.
The tangent direction of each 1D segment is
\[
  \bm b = \frac{\bm x_1 - \bm x_0}{\lVert \bm x_1 - \bm x_0\rVert},
\]
and the projection of $\bm b$ on the face normal $\bm n$ is denoted by
$\beta = \bm b\!\cdot\!\bm n$.
The pressure $P(A)$ follows a constitutive law (e.g.\ Laplace or elastic model), represented as ,
\begin{equation}
\boxed{%
P(A) = P_0 + E \left[\left(\frac{A}{A_0}\right)^{m} - 1\right],
}
\qquad m = \text{tube law exponent.}
\end{equation}

and $dP(A)=\mathrm{d}P/\mathrm{d}A$ enters the weak formulation.

\section{Discrete space and notation}
Let $\Th$ be a partition of $\Gamma$ into 1D cells $K$. For polynomial degree $k$,
\[
  \Vh = \{ v\in L^2(\Gamma): v|_K \in \mathbb P_k(K)\ \forall K\in\Th \},
\]
implemented by \texttt{FE\_DGQ$<1,\,$spacedim$>$}.  
For $F\in\Fh^\circ$ an interior face with unit normal $\bm n$, we set
\[
  \avg{v}=\tfrac12(v^-+v^+), \qquad \jump{v}=v^- - v^+.
\]
On boundary faces $F\subset\partial\Gamma$ we set $\avg{v}=v$ and $\jump{v}=v$.

 

\section{Lax--Friedrichs numerical flux}
The Lax--Friedrichs (LF) flux is a robust, symmetric upwind flux that
stabilizes nonlinear hyperbolic systems by adding dissipation proportional to the
maximum characteristic speed.  
For a generic flux $\mathbf{F}(W)$ and normal $\bm n$, it is defined as
\begin{equation}
  \widehat{\mathbf{F}}(W^-,W^+) \cdot \bm n
  = \tfrac12 \big[ \mathbf{F}(W^-)\cdot\bm n + \mathbf{F}(W^+)\cdot\bm n \big]
    - \tfrac{\alpha}{2}(W^+ - W^-),
  \label{eq:laxfriedrichs-flux}
\end{equation}
where
\[
  \alpha = \max_{s=\pm}\big|\lambda_{\max}\big(\mathbf{H}(W^s)\big)\big|,
  \qquad
  \mathbf{H} = \frac{\partial \mathbf{F}}{\partial W}
\]
is the flux Jacobian and $\lambda_{\max}$ its largest eigenvalue.
For the blood flow system,
\[
  \mathbf{F}(W) =
  \begin{pmatrix}
    A U \\[2mm]
    \dfrac{U^2}{2} + \dfrac{1}{\rho}P(A)
  \end{pmatrix}, \qquad
  \mathbf{H} =
  \begin{pmatrix}
    U & A \\[1mm]
    \dfrac{c^2}{A} & U
  \end{pmatrix}, \qquad
  \lambda_{1,2} = U \pm c, \quad
  c=\sqrt{\dfrac{A}{\rho}\dfrac{dP}{dA}}.
\]
Thus $\alpha = \max(|U|+c)$ on each face.

\section{DG weak formulation with Lax--Friedrichs flux}

We seek $(A_h,U_h)\in\Vh\times\Vh$ such that for all test functions
$(\phi_A,\phi_U)\in\Vh\times\Vh$:
\begin{align}
\sum_{K\in\Th}\!\!\int_K
  \frac{A_h^{n+1}-A_h^n}{\dt}\,\phi_A
&\;-\;
  (A_h^{n+1}\,\bar U_h)\,\bm b\!\cdot\!\nabla\phi_A
\nonumber\\
&\;+\;
  \sum_{F\in\Fh}\!\!\int_F
  \widehat{(A U)}^{\text{LF}}\,\phi_A^-\,\mathrm d\ell
\;=\;
  \sum_{K\in\Th}\!\!\int_K f_A\,\phi_A ,
\label{eq:area-dg-lf}
\\[8pt]
\sum_{K\in\Th}\!\!\int_K
  \frac{U_h^{n+1}-U_h^n}{\dt}\,\phi_U
&\;-\;
  \Big(
    \tfrac{\bar U_h}{2}\,U_h^{n+1}
    + \tfrac{1}{\rho}P(\bar A_h)
  \Big)
  \bm b\!\cdot\!\nabla\phi_U
\nonumber\\
&\;+\;
  \sum_{F\in\Fh}\!\!\int_F
  \widehat{\!\left(
    \tfrac{U^2}{2}
    + \tfrac{1}{\rho}P(A)
  \right)\!}^{\text{LF}}
  \phi_U^-\,\mathrm d\ell
\;=\;
  \sum_{K\in\Th}\!\!\int_K
  (f_U - c\,U_h^{n+1})\,\phi_U .
\label{eq:momentum-dg-lf}
\end{align}

The LF fluxes are computed from \eqref{eq:laxfriedrichs-flux} for each
variable with $\alpha = \max(|U|+c)$ across the face.

 

\section{Boundary conditions and Lax--Friedrichs numerical fluxes}

At boundary faces $F\subset\partial\Gamma$, an exterior (ghost) state
$W^+ = (A^+,U^+)^\top$ is introduced to construct the Lax--Friedrichs (LF)
numerical flux. The generic form of the LF flux for a conservation law
$\partial_t W + \bm b\!\cdot\!\nabla \mathbf{F}(W) = 0$
with normal direction $\bm n$ and $\beta=\bm b\!\cdot\!\bm n$
is given by
\begin{equation}
  \widehat{\mathbf{F}}^{\text{LF}}(W^-,W^+)\cdot\bm n
  = \tfrac12\beta\big[\mathbf{F}(W^-)+\mathbf{F}(W^+)\big]
    - \tfrac{\alpha}{2}(W^+ - W^-),
  \label{eq:LF-flux-general}
\end{equation}
where $\alpha$ is a local dissipation coefficient chosen as
\[
  \alpha = \max\big(|U^-|+c^-,\,|U^+|+c^+\big),
  \qquad
  c = \sqrt{\tfrac{A}{\rho}\,\tfrac{dP}{dA}} .
\]
The sign of the eigenvalues $\lambda_{1,2} = U \pm c$ of the flux Jacobian
$\mathbf{H}=\partial\mathbf{F}/\partial W$
determines the direction of information propagation:
\[
W^+ =
\begin{cases}
  W_{\text{ext}} & \text{if inflow, i.e. } \lambda_i(W^-)\!<\!0,\\[2pt]
  W^- & \text{if outflow, i.e. } \lambda_i(W^-)\!>\!0.
\end{cases}
\]

\subsection*{Lax--Friedrichs flux for the area equation}
For the first component of $\mathbf{F}(W)$, namely
\[
  F_1(A,U) = A\,U ,
\]
the Lax--Friedrichs flux across a face $F$ is
\begin{equation}
  \widehat{(A U)}_{\text{LF}}
  = \tfrac12\beta\big[(A^-U^-)+(A^+U^+)\big]
   - \tfrac{\alpha}{2}(A^+ - A^-),
  \label{eq:LF-flux-AU}
\end{equation}
where $\beta=\bm b\!\cdot\!\bm n$ and $\alpha$ is as defined above.
The first term provides the centered flux, while the second introduces
dissipation proportional to the jump $\jump{A}$.

\subsection*{Lax--Friedrichs flux for the momentum equation}
The second component of the physical flux is
\[
  F_2(A,U) = \frac{U^2}{2} + \frac{1}{\rho}P(A) ,
\]
representing the convective and pressure contributions to the momentum balance.
The corresponding Lax--Friedrichs flux reads
\begin{equation}
  \widehat{\left(\frac{U^2}{2} + \frac{1}{\rho}P(A)\right)}_{\!\!\text{LF}}
  = \tfrac12\beta
    \left[
      \left(\frac{(U^-)^2}{2}+\frac{1}{\rho}P(A^-)\right)
      + \left(\frac{(U^+)^2}{2}+\frac{1}{\rho}P(A^+)\right)
    \right]
   - \tfrac{\alpha}{2}(U^+ - U^-),
  \label{eq:LF-flux-U2PA}
\end{equation}
where again $\alpha=\max(|U|+c)$.

\subsection*{Inflow case}
When $\lambda_1 < 0 \text{ and } \lambda_2 <0$ then:
\begin{itemize}
  \item The LF fluxes \eqref{eq:LF-flux-AU} is equivalent to $\tfrac12\beta (A^-U^-)$ and
  \eqref{eq:LF-flux-U2PA} is equivalent to $\tfrac12\beta
 \left(\frac{(U^-)^2}{2}+\frac{1}{\rho}P(A^-)\right)$.
\end{itemize}
\subsection*{Outflow case}
When $\lambda_1 > 0 \text{ and } \lambda_2 > 0$ then:
      \begin{itemize}
          \item The LF fluxes \eqref{eq:LF-flux-AU} is equivalent to $\tfrac12\beta (A_{bc} U_{bc})$ and
  \eqref{eq:LF-flux-U2PA} is equivalent to $\tfrac12\beta
 \left(\frac{(U_{bc})^2}{2}+\frac{1}{\rho}P(A_{bc})\right)$.
      \end{itemize}

\section{Algorithm summary with Lax--Friedrichs flux}
\begin{enumerate}
  \item Assemble $M_A$, $M_U$.
  \item For each time step:
  \begin{enumerate}
    \item Initialize Picard iterate $(A^{(0)},U^{(0)})=(A^n,U^n)$.
    \item Assemble fluxes using the LF formula \eqref{eq:laxfriedrichs-flux}.
    \item Solve the linearized system with updated $\mathcal{A}$.
    \item Relax and update until Picard convergence.
  \end{enumerate}
  \item Advance to $t^{n+1}$ and write output.
\end{enumerate}

\end{document}
